\chapter{Writeback Stage}

\section{Overview}

The Writeback (WB) stage is responsible for collecting the completed results produced by the
execution units and making them available for architectural retirement. Unlike an in-order
processor where results are written directly into the register file, Sargantana employs an
out-of-order backend. Consequently, completed execution results are first gathered and
associated with their corresponding physical registers before the instruction becomes eligible for
retirement.

The writeback stage therefore acts as the interface between execution and commit. It receives
completion notifications from multiple functional units such as the ALU, multiplier, divider,
floating-point unit, SIMD unit, and memory subsystem. Once a result is available, the stage
updates the internal completion status so that dependent instructions may proceed and the
commit stage can eventually retire the instruction.

Unlike the execution stage, the writeback logic mainly manages completed instruction
information rather than performing arithmetic computations.

\begin{rtlrefs}
  \rtlitem{rtl/datapath/rtl/wb_stage/rtl/graduation_list.sv}
  \rtlitem{rtl/datapath/rtl/datapath.sv}
\end{rtlrefs}

\begin{figure}[H]
  \centering
  \includegraphics[width=\textwidth]{ch11_writeback_overview.png}
  \caption{Writeback stage in context: results flow from the execution units into the Graduation List and Commit}
  \label{fig:writeback-overview}
\end{figure}

\section{Result Collection}

Multiple execution units operate simultaneously and may finish execution at different clock
cycles. The writeback stage therefore collects completed results from all available execution
pipelines.

Typical producers include:
\begin{itemize}
  \item Integer ALU
  \item Branch Unit
  \item Multiplier
  \item Divider
  \item Memory Unit
  \item Floating Point Unit
  \item SIMD Unit
\end{itemize}

\begin{figure}[H]
  \centering
  \includegraphics[width=\textwidth]{ch11_result_collection.png}
  \caption{Result collection: completed results from the highlighted execution units feed the writeback stage}
  \label{fig:result-collection}
\end{figure}

Each completed instruction provides:
\begin{itemize}
  \item Destination physical register
  \item Computed result
  \item Completion indication
  \item Exception information (if generated)
\end{itemize}

Rather than immediately committing the instruction, the writeback stage records that the
instruction has completed successfully.

The RTL implements this functionality through the \textbf{Graduation List}, which keeps track of
instructions that have completed execution and are waiting for architectural retirement.

\begin{rtlrefs}
  \rtlitem{rtl/datapath/rtl/wb_stage/rtl/graduation_list.sv}
\end{rtlrefs}

\section{Register Update}

Because Sargantana uses physical register renaming, execution results are written into the
allocated physical register instead of directly updating the architectural register file.

The sequence is:
\begin{enumerate}
  \item Register Rename allocates a physical register.
  \item Execution computes the instruction result.
  \item Writeback stores the value into the physical register.
  \item Commit later updates the architectural mapping.
\end{enumerate}

This separation allows speculative instructions to execute safely without modifying the visible
architectural state.

The Graduation List maintains the completion status for each instruction until the Commit stage
verifies that the instruction can safely retire.

This mechanism supports:
\begin{itemize}
  \item speculative execution
  \item precise exceptions
  \item rollback after branch misprediction
  \item correct in-order retirement
\end{itemize}

\begin{rtlrefs}
  \rtlitem{rtl/datapath/rtl/ir_stage/rtl/rename_table.sv}
  \rtlitem{rtl/datapath/rtl/rr_stage/rtl/regfile.sv}
  \rtlitem{rtl/datapath/rtl/wb_stage/rtl/graduation_list.sv}
\end{rtlrefs}

\section{Exception Information}

Besides carrying execution results, the writeback stage also forwards exception information
generated during execution.

Examples include:
\begin{itemize}
  \item Illegal instruction
  \item Load or Store access fault
  \item Misaligned memory access
  \item Arithmetic exceptions
  \item Floating-point exceptions
\end{itemize}

Instead of immediately servicing an exception, the writeback stage records the exception status
alongside the completed instruction.

The Commit stage later checks whether the faulting instruction is the oldest instruction in the
machine. Only then is the architectural state updated and the exception is delivered,
guaranteeing precise exception handling.

This approach ensures that younger speculative instructions never modify architectural state
before older instructions have safely completed.

\begin{rtlrefs}
  \rtlitem{rtl/datapath/rtl/wb_stage/rtl/graduation_list.sv}
  \rtlitem{rtl/control_unit/rtl/control_unit.sv}
\end{rtlrefs}
